\section{Finite-Field Consequences}
\label{sec:finite-fields}

Now let \(V=\Fq^n\). The weak-congruence row labels from
Section~\ref{sec:geometric-stabilizers} remain in force, but over \(\Fq\) some
geometric rows split according to the Frobenius action on the underlying simple
support factors. The resulting finite-field stabilizers are recorded in the
tables below.

\begin{proposition}
\label{prop:fixed-points-rewrite}
Let \(L\) be one of the representatives defined over \(\Fq\), and let
\[
  H_L=\Stab_{\GL_n(\overline{\Fq})}(L).
\]
Then
\[
  \Stab_{\GL_n(q)}(L)=H_L^F,
\]
where \(F\) is the standard Frobenius endomorphism.
\end{proposition}

\begin{proof}
Since \(L\) is \(F\)-stable, the stabilizer over \(\Fq\) is the fixed-point
group of the algebraic stabilizer. Indeed,
\[
  g\in \GL_n(q)
  \iff
  F(g)=g.
\]
If such a \(g\) stabilizes \(L\), then it lies in \(H_L^F\). Conversely every
element of \(H_L^F\) is an \(F\)-fixed element of \(\GL_n(\overline{\Fq})\) that
stabilizes \(L\), hence belongs to \(\GL_n(q)\) and stabilizes \(L\) over
\(\Fq\). Therefore
\[
  \Stab_{\GL_n(q)}(L)=H_L^F.
\]
\end{proof}

\begin{proposition}
\label{prop:finite-stabilizers-rewrite}
Every finite-field stabilizer through \(n\le 7\) is obtained from the
corresponding algebraic stabilizer by Frobenius fixed points. Distinct support
points in one Frobenius orbit give the field-extension groups
\(\SL_2(q^2)\) and \(\SL_2(q^3)\). Repeated support at one rational point gives
the groups \(U_{3d-3}(q)\rtimes \SL_2(q)\).

The orbit-by-orbit formulas are listed in the finite-field tables below.
\end{proposition}

\begin{proof}
Let \(L\) be an \(\Fq\)-rational representative. Proposition
\ref{prop:fixed-points-rewrite} identifies its finite-field stabilizer with the
fixed-point group of the algebraic stabilizer. We therefore only need to
understand how Frobenius acts on the quotient \(Q_L\) and on the kernel \(K_L\).

On the quotient side, rational support gives the split families already visible
over \(k\). If two or three simple-support factors form one Frobenius orbit,
then the induced permutation on the coefficient line produces the non-split
torus normalizer or the cyclic quadratic/cubic quotient recorded in the
finite-field tables.

On the kernel side, there are two cases. If the support points are distinct and
lie in one Frobenius orbit, the corresponding simple \(\SL_2(k)\)-factors are
identified by Frobenius, and the fixed points are the field-extension groups
\(\SL_2(q^2)\) or \(\SL_2(q^3)\). If one support point is repeated, then the
intersection calculation already lives over the rational point, so Frobenius
acts coefficientwise on the parameters appearing in the unipotent radical and
on the Levi factor \(\SL_2(k)\). The fixed-point group therefore has
unipotent radical of dimension \(3d-3\) and Levi quotient \(\SL_2(q)\).
Therefore it is written in the tables as
\(U_{3d-3}(q)\rtimes \SL_2(q)\). Radical extensions descend termwise because
their defining formulas are polynomial in the coefficients of the
representative. The same is true for the one-flat extensions from the
geometric table.
\end{proof}

\medskip

Here \(\lambda,\mu,\nu\in \Fq\) denote distinct rational parameters. The symbol
\(Q\) denotes the quadratic simple-support row, and \(C\) denotes the cubic
simple-support row. In the quotient column, \(B(q)\), \(T_s(q)\), \(N_s(q)\),
and \(N_{ns}(q)\) denote the split Borel subgroup of \(\GL_2(q)\), the split
maximal torus, the normalizer of the split torus, and the normalizer of a
non-split torus. The groups \(\widetilde{S}_3(q)\), \(\widetilde{C}_2(q)\),
and \(\widetilde{C}_3(q)\) are the corresponding full preimages in \(\GL_2(q)\)
of \(S_3\), \(C_2\), and \(C_3\) in \(\PGL_2(q)\).

\begingroup
\footnotesize
\setlength{\tabcolsep}{4pt}
% This file is generated by magma/examples/generate_geometric_finite_field_stabilizer_tables_tex.m.

\paragraph{$n=4$.}
\begin{longtable}{@{}c p{0.28\textwidth} p{0.47\textwidth} p{0.12\textwidth}@{}}
\toprule
No. & Type & $K_L$ in $\Stab_{\GL_n(q)}(L)=K_L\rtimes Q_L$ & $Q_L$ \\
\midrule
\endfirsthead
\toprule
No. & Type & $K_L$ in $\Stab_{\GL_n(q)}(L)=K_L\rtimes Q_L$ & $Q_L$ \\
\midrule
\endhead
\midrule
\multicolumn{4}{r}{Continued on next page}\\
\midrule
\endfoot
\bottomrule
\endlastfoot
1 & \(J_{\lambda,2}\) & $U_3(q)\rtimes \SL_2(q)$ & $B(q)$ \\
2 & \(J_{\lambda,1} + J_{\mu,1}\ (\lambda \ne \mu)\) & $\SL_2(q)\times \SL_2(q)$ & $N_s(q)$ \\
3 & \(Q\) & $\SL_2(q^2)$ & $N_{ns}(q)$ \\
4 & \(S_1 + S_2\) & $U_5(q)\rtimes (\GL_1(q)\times \GL_1(q))$ & $\GL_2(q)$ \\
\end{longtable}

\paragraph{$n=5$.}
\begin{longtable}{@{}c p{0.28\textwidth} p{0.47\textwidth} p{0.12\textwidth}@{}}
\toprule
No. & Type & $K_L$ in $\Stab_{\GL_n(q)}(L)=K_L\rtimes Q_L$ & $Q_L$ \\
\midrule
\endfirsthead
\toprule
No. & Type & $K_L$ in $\Stab_{\GL_n(q)}(L)=K_L\rtimes Q_L$ & $Q_L$ \\
\midrule
\endhead
\midrule
\multicolumn{4}{r}{Continued on next page}\\
\midrule
\endfoot
\bottomrule
\endlastfoot
1 & \(S_3\) & $U_4(q)\rtimes \GL_1(q)$ & $\GL_2(q)$ \\
2 & \(S_2 + S_1^2\) & $U_8(q)\rtimes (\GL_2(q)\times \GL_1(q))$ & $\GL_2(q)$ \\
3 & \(S_2 + J_{\lambda,1}\) & $U_4(q)\rtimes (\GL_1(q)\times \SL_2(q))$ & $B(q)$ \\
4 & \(S_1 + J_{\lambda,2}\) & $U_7(q)\rtimes (\GL_1(q)\times \SL_2(q))$ & $B(q)$ \\
5 & \(S_1 + J_{\lambda,1} + J_{\mu,1}\ (\lambda \ne \mu)\) & $U_4(q)\rtimes (\GL_1(q)\times \SL_2(q)\times \SL_2(q))$ & $N_s(q)$ \\
6 & \(S_1 + Q\) & $U_4(q)\rtimes (\GL_1(q)\times \SL_2(q^2))$ & $N_{ns}(q)$ \\
\end{longtable}

\paragraph{$n=6$.}
\begin{longtable}{@{}c p{0.28\textwidth} p{0.47\textwidth} p{0.12\textwidth}@{}}
\toprule
No. & Type & $K_L$ in $\Stab_{\GL_n(q)}(L)=K_L\rtimes Q_L$ & $Q_L$ \\
\midrule
\endfirsthead
\toprule
No. & Type & $K_L$ in $\Stab_{\GL_n(q)}(L)=K_L\rtimes Q_L$ & $Q_L$ \\
\midrule
\endhead
\midrule
\multicolumn{4}{r}{Continued on next page}\\
\midrule
\endfoot
\bottomrule
\endlastfoot
1 & \(J_{\lambda,3}\) & $U_6(q)\rtimes \SL_2(q)$ & $B(q)$ \\
2 & \(J_{\lambda,2} + J_{\lambda,1}\) & $U_7(q)\rtimes (\SL_2(q)\times \SL_2(q))$ & $B(q)$ \\
3 & \(J_{\lambda,2} + J_{\mu,1}\ (\lambda \ne \mu)\) & $U_3(q)\rtimes (\SL_2(q)\times \SL_2(q))$ & $T_s(q)$ \\
4 & \(2J_{\lambda,1} + J_{\mu,1}\ (\lambda \ne \mu)\) & $\Sp_4(q)\times \SL_2(q)$ & $T_s(q)$ \\
5 & \(J_{\lambda,1} + J_{\mu,1} + J_{\nu,1}\ \text{(distinct)}\) & $\SL_2(q)\times \SL_2(q)\times \SL_2(q)$ & $\widetilde{S}_3(q)$ \\
6 & \(J_{\lambda,1} + Q\) & $\SL_2(q)\times \SL_2(q^2)$ & $\widetilde{C}_2(q)$ \\
7 & \(C\) & $\SL_2(q^3)$ & $\widetilde{C}_3(q)$ \\
8 & \(S_1^2 + J_{\lambda,2}\) & $U_{11}(q)\rtimes (\GL_2(q)\times \SL_2(q))$ & $B(q)$ \\
9 & \(S_1^2 + J_{\lambda,1} + J_{\mu,1}\ (\lambda \ne \mu)\) & $U_8(q)\rtimes (\GL_2(q)\times \SL_2(q)\times \SL_2(q))$ & $N_s(q)$ \\
10 & \(S_1^2 + Q\) & $U_8(q)\rtimes (\GL_2(q)\times \SL_2(q^2))$ & $N_{ns}(q)$ \\
11 & \(S_1 + S_2 + J_{\lambda,1}\) & $U_9(q)\rtimes (\GL_1(q)\times \GL_1(q)\times \SL_2(q))$ & $B(q)$ \\
12 & \(S_1^3 + S_2\) & $U_{11}(q)\rtimes (\GL_3(q)\times \GL_1(q))$ & $\GL_2(q)$ \\
13 & \(S_1 + S_3\) & $U_8(q)\rtimes (\GL_1(q)\times \GL_1(q))$ & $\GL_2(q)$ \\
14 & \(S_2^2\) & $U_6(q)\rtimes \GL_2(q)$ & $\GL_2(q)$ \\
\end{longtable}

\paragraph{$n=7$.}
\begin{longtable}{@{}c p{0.28\textwidth} p{0.47\textwidth} p{0.12\textwidth}@{}}
\toprule
No. & Type & $K_L$ in $\Stab_{\GL_n(q)}(L)=K_L\rtimes Q_L$ & $Q_L$ \\
\midrule
\endfirsthead
\toprule
No. & Type & $K_L$ in $\Stab_{\GL_n(q)}(L)=K_L\rtimes Q_L$ & $Q_L$ \\
\midrule
\endhead
\midrule
\multicolumn{4}{r}{Continued on next page}\\
\midrule
\endfoot
\bottomrule
\endlastfoot
1 & \(S_4\) & $U_6(q)\rtimes \GL_1(q)$ & $\GL_2(q)$ \\
2 & \(S_3 + S_1^2\) & $U_{14}(q)\rtimes (\GL_2(q)\times \GL_1(q))$ & $\GL_2(q)$ \\
3 & \(S_2^2 + S_1\) & $U_{12}(q)\rtimes (\GL_1(q)\times \GL_2(q))$ & $\GL_2(q)$ \\
4 & \(S_2 + S_1^4\) & $U_{14}(q)\rtimes (\GL_4(q)\times \GL_1(q))$ & $\GL_2(q)$ \\
5 & \(S_3 + J_{\lambda,1}\) & $U_6(q)\rtimes (\GL_1(q)\times \SL_2(q))$ & $B(q)$ \\
6 & \(S_2 + S_1^2 + J_{\lambda,1}\) & $U_{14}(q)\rtimes (\GL_2(q)\times \GL_1(q)\times \SL_2(q))$ & $B(q)$ \\
7 & \(S_2 + J_{\lambda,2}\) & $U_9(q)\rtimes (\GL_1(q)\times \SL_2(q))$ & $B(q)$ \\
8 & \(S_2 + 2J_{\lambda,1}\) & $U_6(q)\rtimes (\GL_1(q)\times \Sp_4(q))$ & $B(q)$ \\
9 & \(S_2 + J_{\lambda,1} + J_{\mu,1}\ (\lambda \ne \mu)\) & $U_6(q)\rtimes (\GL_1(q)\times \SL_2(q)\times \SL_2(q))$ & $N_s(q)$ \\
10 & \(S_2 + Q\) & $U_6(q)\rtimes (\GL_1(q)\times \SL_2(q^2))$ & $N_{ns}(q)$ \\
11 & \(S_1^3 + J_{\lambda,2}\) & $U_{15}(q)\rtimes (\GL_3(q)\times \SL_2(q))$ & $B(q)$ \\
12 & \(S_1^3 + J_{\lambda,1} + J_{\mu,1}\ (\lambda \ne \mu)\) & $U_{12}(q)\rtimes (\GL_3(q)\times \SL_2(q)\times \SL_2(q))$ & $N_s(q)$ \\
13 & \(S_1^3 + Q\) & $U_{12}(q)\rtimes (\GL_3(q)\times \SL_2(q^2))$ & $N_{ns}(q)$ \\
14 & \(S_1 + J_{\lambda,3}\) & $U_{12}(q)\rtimes (\GL_1(q)\times \SL_2(q))$ & $B(q)$ \\
15 & \(S_1 + J_{\lambda,2} + J_{\lambda,1}\) & $U_{13}(q)\rtimes (\GL_1(q)\times \SL_2(q)\times \SL_2(q))$ & $B(q)$ \\
16 & \(S_1 + J_{\lambda,2} + J_{\mu,1}\ (\lambda \ne \mu)\) & $U_9(q)\rtimes (\GL_1(q)\times \SL_2(q)\times \SL_2(q))$ & $T_s(q)$ \\
17 & \(S_1 + 2J_{\lambda,1} + J_{\mu,1}\ (\lambda \ne \mu)\) & $U_6(q)\rtimes (\GL_1(q)\times \Sp_4(q)\times \SL_2(q))$ & $T_s(q)$ \\
18 & \(S_1 + J_{\lambda,1} + J_{\mu,1} + J_{\nu,1}\ \text{(distinct)}\) & $U_6(q)\rtimes (\GL_1(q)\times \SL_2(q)\times \SL_2(q)\times \SL_2(q))$ & $\widetilde{S}_3(q)$ \\
19 & \(S_1 + J_{\lambda,1} + Q\) & $U_6(q)\rtimes (\GL_1(q)\times \SL_2(q)\times \SL_2(q^2))$ & $\widetilde{C}_2(q)$ \\
20 & \(S_1 + C\) & $U_6(q)\rtimes (\GL_1(q)\times \SL_2(q^3))$ & $\widetilde{C}_3(q)$ \\
\end{longtable}

\endgroup

\begin{theorem}
\label{thm:orbit-stabilizer-completeness}
For each \(n=4,5,6,7\), let \(\mathcal T_n(q)\) be the finite list of orbit
representatives in the finite-field tables above. Then
\[
  \sum_{L\in \mathcal T_n(q)}
  \frac{|\GL_n(q)|}{|\Stab_{\GL_n(q)}(L)|}
  =
  \left[\begin{matrix}\binom{n}{2}\\2\end{matrix}\right]_q.
\]
Therefore the finite-field tables above are the complete orbit table for
\(\GL_n(q)\) on \(\Gr(2,\wedge^2 \Fq^n)\).
\end{theorem}

\begin{proof}
The right-hand side is the number of \(2\)-spaces in \(\wedge^2 \Fq^n\). The
left-hand side is the total number of elements covered by the listed orbits.
For each row of the finite-field tables, the orbit-stabilizer theorem gives an
orbit of size
\[
  \frac{|\GL_n(q)|}{|\Stab_{\GL_n(q)}(L)|}.
\]
Summing over all listed representatives therefore counts the total number of
points lying in the listed orbits.

It remains to compare this sum with the number of all \(2\)-spaces in
\(\wedge^2\Fq^n\), which is the Gaussian polynomial on the right. Substituting
the explicit stabilizer formulas from the finite-field tables gives a rational
function in \(q\). After clearing denominators, one obtains a polynomial
identity. The symbolic simplification is routine but lengthy, and we verify it
with the companion Magma scripts archived with the paper. Therefore the listed
orbit sizes add up to the whole Grassmannian.

Since distinct orbits are disjoint, no further orbit can exist. Therefore the
finite-field tables are complete.
\end{proof}

\begin{corollary}
\label{cor:finite-counts-rewrite}
The numbers of \(\GL_n(q)\)-orbits on \(\Gr(2,\wedge^2 \Fq^n)\) are
\[
  4,\ 6,\ 14,\ 20
\]
for \(n=4,5,6,7\), respectively.
\end{corollary}

\begin{proof}
The finite-field tables above have exactly \(4\), \(6\), \(14\), and \(20\)
rows for \(n=4\), \(5\), \(6\), and \(7\), respectively.
\end{proof}

\begin{corollary}
\label{cor:geometric-list-complete-rewrite}
The geometric orbit representatives in
Section~\ref{sec:geometric-stabilizers} form a complete set of
\(\GL(V)\)-orbit representatives on \(\Gr(2,\wedge^2 V)\) for \(\dim V=4,5,6,7\).
The numbers of geometric orbits are
\[
  3,\ 5,\ 11,\ 16
\]
for \(\dim V=4,5,6,7\), respectively.
\end{corollary}

\begin{proof}
Let \(k=\overline{\Fq}\), let \(V=k^n\) with \(n\le 7\), and let
\(L\le \wedge^2 V\) be any geometric \(2\)-space. Choose a basis
\(\omega_1,\omega_2\) of \(L\). Since \(k\) is the algebraic closure of \(\Fq\),
the finitely many coefficients of \(\omega_1\) and \(\omega_2\) lie in some
finite extension \(\mathbf F_{q^m}\subset k\). Therefore \(L\) is defined over
\(\mathbf F_{q^m}\).

Apply Theorem~\ref{thm:orbit-stabilizer-completeness} with \(q\) replaced by
\(q^m\). The resulting finite-field orbit table is complete for
\(\GL_n(q^m)\) on \(\Gr(2,\wedge^2 \mathbf F_{q^m}^n)\). Hence \(L\) is
\(\GL_n(q^m)\)-equivalent to one of the finite-field representatives listed in
this section. Each of those finite-field representatives is obtained by
Frobenius descent from one of the geometric representatives in
Section~\ref{sec:geometric-stabilizers}. After extending scalars back to \(k\),
the same equivalence holds in \(\GL_n(k)\). Therefore every geometric orbit is
represented by one of the listed geometric rows.

Counting the rows in Section~\ref{sec:geometric-stabilizers} gives \(3\), \(5\),
\(11\), and \(16\) representatives in dimensions \(4\), \(5\), \(6\), and
\(7\), respectively.
\end{proof}
